\documentclass[border=18pt]{standalone}

% --- Paquetes de Codificación y Tipografía ---
\usepackage[utf8]{inputenc}
\usepackage[T1]{fontenc}
\usepackage{lmodern}

% --- Paquetes Matemáticos y Físicos ---
\usepackage{amsmath}
\usepackage{amssymb}
\usepackage{siunitx}

% --- Monocromatismo Estricto (Forzar Blanco y Negro 100%) ---
\usepackage{xcolor}
\definecolor{black}{HTML}{000000}
\definecolor{white}{HTML}{FFFFFF}

% --- Núcleo de Diagramación: TikZ y CircuiTikZ ---
\usepackage{tikz}
\usepackage[american, cute inductors, american currents, american voltages]{circuitikz}

% --- Librerías Especializadas TikZ ---
\usetikzlibrary{calc}
\usetikzlibrary{positioning}
\usetikzlibrary{backgrounds}

% --- Declaración y Configuración de Capas PGF ---
\pgfdeclarelayer{background}
\pgfdeclarelayer{foreground}
\pgfsetlayers{background,main,foreground}

% --- Configuración Estándar IEEE para CircuiTikZ ---
\ctikzset{
    logic ports=ieee,
    resistors/scale=0.70,
    capacitors/scale=0.65,
    diodes/scale=0.65,
    transistors/scale=0.80,
    sources/scale=0.70,
    bipoles/thickness=1.0,
    nodes width=0.055
}

% --- Definición de Estilos Personalizados 100% Monocromáticos ---
\tikzset{
    rail/.style={line width=1.1pt, draw=black},
    signal/.style={line width=0.75pt, draw=black},
    feedback_bus/.style={line width=0.80pt, draw=black},
    probe_line/.style={draw=gray!70, densely dashed, line width=0.75pt},
    box_frame/.style={draw=gray!50, dashed, rounded corners=4pt, line width=0.6pt},
    scope_frame/.style={draw=black, fill=white, rounded corners=3pt, line width=0.85pt},
    scope_grid/.style={step=0.25cm, gray!25, very thin},
    scope_axis/.style={gray!70, thin},
    scope_wave/.style={draw=black, line width=1.15pt}
}

\begin{document}
\begin{circuitikz}[line width=0.75pt, american]

% =========================================================================
% 🗺️ CUADRÍCULA MAESTRA GLOBAL (Master Coordinate Grid)
% Riel Superior (V_RAW):          Y = 9.0
% Riel de Salida (V_OUT):         Y = 6.73
% Riel Inferior (GND):            Y = 0.0
% Lazo de Realimentación (FB):    Y = 10.0
% Capa de Instrumentación:        Y = 11.5 a 13.8
% =========================================================================

% -------------------------------------------------------------------------
% 📊 CAPA ANALÍTICA: 5 OSCILOSCOPIOS MATEMÁTICOS (Background Layer)
% Puras representaciones visuales sin títulos ni texto descriptivo.
% -------------------------------------------------------------------------
\begin{pgfonlayer}{background}

    % Macro Interna de Osciloscopio (Silencio Textual Estricto)
    \newcommand{\renderScope}[3]{
        \begin{scope}[shift={(#1,#2)}]
            % Chasis del Instrumento
            \draw[scope_frame] (-1.25, -0.95) rectangle (1.25, 0.95);
            % Graticule (Pantalla Cuadriculada)
            \draw[scope_grid] (-1.1, -0.8) grid (1.1, 0.8);
            % Ejes de Referencia
            \draw[scope_axis] (-1.1, 0) -- (1.1, 0);
            \draw[scope_axis] (0, -0.8) -- (0, 0.8);
            % Trazo Matemático de Señal
            \begin{scope}
                \clip (-1.1, -0.8) rectangle (1.1, 0.8);
                #3
            \end{scope}
        \end{scope}
    }

    % Osciloscopio 1: Red AC (Onda Senoidal Pura)
    \renderScope{1.0}{12.5}{
        \draw[scope_wave] plot[domain=-1.1:1.1, samples=60] (\x, {0.65*sin(360*\x/1.1)});
    }

    % Osciloscopio 2: Secundario Transformador (Seno Reducido y Aislado)
    \renderScope{4.8}{12.5}{
        \draw[scope_wave] plot[domain=-1.1:1.1, samples=60] (\x, {0.45*sin(360*\x/1.1)});
    }

    % Osciloscopio 3: Rectificación de Onda Completa (|sin|)
    \renderScope{9.5}{12.5}{
        \draw[scope_wave] plot[domain=-1.1:1.1, samples=80] (\x, {0.70*abs(sin(360*\x/1.1)) - 0.25});
    }

    % Osciloscopio 4: Filtrado Bulk y Rizado Exponencial
    \renderScope{12.2}{12.5}{
        \draw[scope_wave] plot[domain=-1.1:1.1, samples=100] (\x, {0.35 + 0.18*(1 - mod(\x+1.1, 0.55)/0.55)});
    }

    % Osciloscopio 5: Tensión DC Regulada Estable (Línea Plana)
    \renderScope{40.6}{12.5}{
        \draw[scope_wave] (-1.1, 0.35) -- (1.1, 0.35);
    }

    % Sondas de Medición (Test Probes) con Ruteo Ortogonal Estricto y Puntos Rellenos
    % Sonda 1 -> Borne Superior Fuente AC
    \draw[probe_line] (1.0, 11.55) -- (1.0, 4.6);
    \fill[black] (1.0, 4.6) circle (1.8pt);

    % Sonda 2 -> Terminal B1 del Secundario Transformador
    \draw[probe_line] (4.8, 11.55) -- (4.8, 4.6);
    \fill[black] (4.8, 4.6) circle (1.8pt);

    % Sonda 3 -> Cátodos del Puente Rectificador (Salida DC no filtrada)
    \draw[probe_line] (9.5, 11.55) -- (9.5, 9.0);
    \fill[black] (9.5, 9.0) circle (1.8pt);

    % Sonda 4 -> Borne Positivo del Capacitor Bulk
    \draw[probe_line] (12.2, 11.55) -- (12.2, 9.0);
    \fill[black] (12.2, 9.0) circle (1.8pt);

    % Sonda 5 -> Riel de Salida Regulada
    \draw[probe_line] (40.6, 11.55) -- (40.6, 6.73);
    \fill[black] (40.6, 6.73) circle (1.8pt);

\end{pgfonlayer}


% =========================================================================
% ⚡ RIELES PRINCIPALES DE ALIMENTACIÓN Y TIERRA
% =========================================================================
% Riel de Tierra Principal (GND) continuo
\draw[rail] (0,0) -- (41.8,0);
\draw (41.4,0) node[ground] {};
\node[font=\sffamily\bfseries\scriptsize, below=4pt] at (2.0, 0.0) {$\text{GND}$};

% Riel de Alimentación no Regulada (V_RAW)
\draw[rail] (9.5,9.0) -- (31.48,9.0) node[right, font=\sffamily\bfseries\scriptsize] {$V_{\text{RAW}}$};


% =========================================================================
% 🏗️ NIVEL 1: MACRO / POTENCIA (Transformador, Puente y Filtro)
% =========================================================================

% --- Transformador con Núcleo de Hierro (Aislamiento Galvánico) ---
\node[transformer core] (T) at (3.0,4.0) {};
\node[font=\sffamily\scriptsize\bfseries, below=2pt] at (T.base) {$T_1$};

% Fuente de Entrada AC (Alineación Horizontal Exacta con Bornes Primarios)
\draw (1.0, 0 |- T.A2) to[sV, l=$v_{\text{AC}}$] (1.0, 0 |- T.A1);
\draw[signal] (1.0, 0 |- T.A1) -- (T.A1);
\draw[signal] (1.0, 0 |- T.A2) -- (T.A2);

% Terminales Secundario (B1, B2)
\coordinate (VAC1_NODE) at (4.8, 0 |- T.B1);
\coordinate (VAC2_NODE) at (4.8, 0 |- T.B2);
\draw[signal] (T.B1) -- (VAC1_NODE);
\draw[signal] (T.B2) -- (VAC2_NODE);
\node[circ, fill=black, inner sep=1.2pt] at (VAC1_NODE) {};
\node[circ, fill=black, inner sep=1.2pt] at (VAC2_NODE) {};

% --- Puente Rectificador en Diamante Planar (100% Libre de Cruces) ---
% Vértices: Superior (AC1), Inferior (AC2), Izquierdo (GND), Derecho (VRAW)
\coordinate (BR_AC1) at (7.8, 5.8);
\coordinate (BR_AC2) at (7.8, 2.2);
\coordinate (BR_GND) at (6.3, 4.0);
\coordinate (BR_VRAW) at (9.3, 4.0);

% Conexión del Secundario al Puente (Ruteo Ortogonal Planar Directo)
\draw[signal] (VAC1_NODE) -- (5.6, 0 |- VAC1_NODE) |- (BR_AC1);
\draw[signal] (VAC2_NODE) -- (5.6, 0 |- VAC2_NODE) |- (BR_AC2);

% 4 Diodos en Diamante Planar
\draw[signal] (BR_AC1) to[D, l=$D_1$] (BR_VRAW);
\draw[signal] (BR_AC2) to[D, l_=$D_2$] (BR_VRAW);
\draw[signal] (BR_GND) to[D, l=$D_3$] (BR_AC1);
\draw[signal] (BR_GND) to[D, l_=$D_4$] (BR_AC2);

\node[circ, fill=black, inner sep=1.2pt] at (BR_AC1) {};
\node[circ, fill=black, inner sep=1.2pt] at (BR_AC2) {};
\node[circ, fill=black, inner sep=1.2pt] at (BR_GND) {};
\node[circ, fill=black, inner sep=1.2pt] at (BR_VRAW) {};

% Conexión Salida DC del Puente a Rieles
\draw[signal] (BR_GND) -- (6.3, 0.0) node[circ, fill=black, inner sep=1.2pt] {};
\draw[signal] (BR_VRAW) -- (9.5, 4.0) -- (9.5, 9.0) node[circ, fill=black, inner sep=1.2pt] {};

% --- Filtro de Entrada (Bulk Electrolítico + Cerámico HF) ---
% Capacitor Electrolítico Bulk (flujo positivo a negativo)
\draw[signal] (12.2, 9.0) to[eC, l=$C_{\text{bulk}}$, *-*] (12.2, 0.0);
% Capacitor Cerámico de Alta Frecuencia
\draw[signal] (13.6, 9.0) to[C, l=$C_{\text{in}}$, *-*] (13.6, 0.0);


% =========================================================================
% 🔬 NIVEL 2: SILICIO / REGULADOR LINEAL
% =========================================================================

% --- Referencia de Voltaje Zener Estable ---
\coordinate (VREF) at (15.2, 4.0);
\draw[signal] (15.2, 9.0) to[R, l=$R_{\text{bias}}$, *-*] (VREF);
\draw[signal] (VREF) to[zD, l_=$D_Z$, -*] (15.2, 0.0);
\node[font=\sffamily\tiny\bfseries, left=1pt] at (VREF) {$V_{\text{ref}}$};


% =========================================================================
% ⚡ NIVEL 3: TRANSISTOR / AMPLIFICADOR DE ERROR DISCRETO
% =========================================================================

% Marco Delimitador de Aislamiento Lógico (Sin texto descriptivo)
\draw[box_frame] (16.0, 0.4) rectangle (30.4, 8.6);

% --- 1. Par Diferencial de Entrada NPN (Q1, Q2) ---
\node[npn] (Q1) at (17.2, 4.0) {};
\node[npn, xscale=-1] (Q2) at (20.2, 4.0) {};
\node[font=\sffamily\tiny\bfseries, above left=1pt] at (Q1.B) {$Q_1$};
\node[font=\sffamily\tiny\bfseries, above right=1pt] at (Q2.B) {$Q_2$};

% Conexión Base Q1 -> Voltaje de Referencia Vref (Horizontal Puro)
\draw[signal] (VREF) -- (Q1.B);

% Conexión Emisores a Fuente de Corriente de Cola
\coordinate (TAIL) at (18.7, 2.7);
\draw[signal] (Q1.E) |- (TAIL) node[circ, fill=black, inner sep=1.2pt] {} -| (Q2.E);
\draw[signal] (TAIL) to[isource, l_=$I_{\text{tail}}$, -*] (18.7, 0.0);

% --- 2. Carga Activa: Espejo de Corriente PNP (Q3, Q4) ---
\node[pnp] (Q3) at (17.2, 6.8) {};
\node[pnp, xscale=-1] (Q4) at (20.2, 6.8) {};
\node[font=\sffamily\tiny\bfseries, below left=2pt] at (Q3.B) {$Q_3$};
\node[font=\sffamily\tiny\bfseries, below right=2pt] at (Q4.B) {$Q_4$};

% Conexión Emisores a V_RAW (Verticales Exactas)
\draw[signal] (Q3.E) -- (Q3.E |- 0, 9.0) node[circ, fill=black, inner sep=1.2pt] {};
\draw[signal] (Q4.E) -- (Q4.E |- 0, 9.0) node[circ, fill=black, inner sep=1.2pt] {};

% Conexión Colectores (Q1 -> Q3, Q2 -> Q4) (Verticales Exactas)
\draw[signal] (Q1.C) -- (Q3.C);
\draw[signal] (Q2.C) -- (Q4.C) coordinate[midway] (N2);

% Interconexión Bases del Espejo de Corriente (Conexión Diodo en Q3)
\draw[signal] (Q3.B) -- (Q4.B);
\draw[signal] (Q3.C) |- (18.7, 6.8) node[circ, fill=black, inner sep=1.2pt] {};
\node[circ, fill=black, inner sep=1.2pt] at (Q3.C) {};

% --- 3. Etapa Amplificadora de Voltaje (VAS) & Compensación Miller ---
\node[npn] (Q5) at (23.2, 4.0) {};
\node[font=\sffamily\tiny\bfseries, above right=1pt] at (Q5.B) {$Q_5$};

% Conexión Salida Diferencial (N2) a Base de Q5 (Ruteo Ortogonal 90°)
\coordinate (VAS_IN) at (21.6, 5.4);
\draw[signal] (N2) node[circ, fill=black, inner sep=1.2pt] {} -- (VAS_IN) |- (Q5.B);

% Emisor de Q5 con Resistencia de Degeneración a Tierra
\draw[signal] (Q5.E) to[R, l_=$R_5$, -*] (Q5.E |- 0, 0.0);

% --- 4. Cadena de Polarización de Drivers (R4, D5, D6) y Acople de Colector VAS ---
% Resistencia Pull-Up R4 desde V_RAW alimentando el ánodo superior (Base Q6)
\coordinate (BIAS_TOP) at (25.2, 7.2);
\draw[signal] (25.2, 9.0) node[circ, fill=black, inner sep=1.2pt] {} to[R, l=$R_4$] (BIAS_TOP) node[circ, fill=black, inner sep=1.2pt] {};

% Cadena de Diodos D5 y D6 en Rama Vertical Estricta
\coordinate (BIAS_BOT) at (25.2, 4.8);
\draw[signal] (BIAS_TOP) to[D, l=$D_5$] (25.2, 6.0) to[D, l=$D_6$] (BIAS_BOT) node[circ, fill=black, inner sep=1.2pt] {};

% Colector de Q5 conectado verticalmente hasta Y=5.4 (nodo Miller) y horizontalmente a BIAS_BOT (Y=4.8)
\coordinate (VAS_COL_TAP) at (Q5.C |- 0, 5.4);
\draw[signal] (Q5.C) -- (VAS_COL_TAP) node[circ, fill=black, inner sep=1.2pt] {};
\draw[signal] (Q5.C |- 0, 4.8) node[circ, fill=black, inner sep=1.2pt] {} -- (BIAS_BOT);

% Capacitor de Compensación Miller (Cc) de Colector de VAS a Entrada de VAS (Horizontal Puro)
\draw[signal] (VAS_COL_TAP) to[C, l_=$C_c$] (VAS_IN);
\node[circ, fill=black, inner sep=1.2pt] at (VAS_IN) {};

% --- 5. Driver Push-Pull Clase AB (Q6, Q7) ---
\node[npn] (Q6) at (27.2, 7.2) {};
\node[pnp] (Q7) at (27.2, 4.8) {};
\node[font=\sffamily\tiny\bfseries, above left=1pt] at (Q6.B) {$Q_6$};
\node[font=\sffamily\tiny\bfseries, below left=1pt] at (Q7.B) {$Q_7$};

% Conexiones de Base a la Cadena de Diodos (Horizontales Puros)
\draw[signal] (BIAS_TOP) -- (Q6.B);
\draw[signal] (BIAS_BOT) -- (Q7.B);

% Colectores del Driver a Rieles de Potencia (Verticales Exactas)
\draw[signal] (Q6.C) -- (Q6.C |- 0, 9.0) node[circ, fill=black, inner sep=1.2pt] {};
\draw[signal] (Q7.C) -- (Q7.C |- 0, 0.0) node[circ, fill=black, inner sep=1.2pt] {};

% Emisores con Degeneración unidos en el Nodo de Excitación (DRV)
\coordinate (DRV) at (29.2, 6.0);
\draw[signal] (Q6.E) to[R, l=$R_{e1}$] (29.2, 6.43) -- (DRV);
\draw[signal] (Q7.E) to[R, l_=$R_{e2}$] (29.2, 5.57) -- (DRV);
\node[circ, fill=black, inner sep=1.2pt] at (DRV) {};
\node[font=\sffamily\tiny\bfseries, above left=1pt] at (DRV) {$\text{DRV}$};


% =========================================================================
% 🛡️ NIVEL 2 (CONTINUACIÓN): PASO DE POTENCIA, SENSADO Y RETROALIMENTACIÓN
% =========================================================================

% --- Transistor de Paso de Potencia (Q_pass NPN) ---
\node[npn] (Qpass) at (31.2, 7.5) {};
\node[font=\sffamily\scriptsize\bfseries, above left=1pt] at (Qpass.B) {$Q_{\text{pass}}$};

% Colector a V_RAW (Vertical Exacta)
\draw[signal] (Qpass.C) -- (Qpass.C |- 0, 9.0) node[circ, fill=black, inner sep=1.2pt] {};
% Base excitada desde el Nodo DRV (Ortogonal 90°)
\draw[signal] (DRV) -- (30.2, 6.0) |- (Qpass.B);

% --- Sensado de Corriente y Protección (R_sc y Q_limit) ---
\coordinate (SENSE) at (32.8, 0 |- Qpass.E);
\draw[signal] (Qpass.E) -- (SENSE) node[circ, fill=black, inner sep=1.2pt] {};

% Resistencia de Sensado en Serie con la Salida
\coordinate (VOUT_NODE) at (35.2, 0 |- Qpass.E);
\draw[signal] (SENSE) to[R, l=$R_{\text{sc}}$] (VOUT_NODE);
\node[circ, fill=black, inner sep=1.2pt] at (VOUT_NODE) {};

% Transistor Limitador de Corriente (Q_limit)
\node[npn] (Qlim) at (31.8, 4.8) {};
\node[font=\sffamily\tiny\bfseries, below left=1pt] at (Qlim.B) {$Q_{\text{limit}}$};

% Sensado Diferencial sobre R_sc (100% Ortogonal)
\draw[signal] (SENSE) -- (32.8, 5.6) -| (Qlim.B);
\draw[signal] (Qlim.E) |- (35.2, 4.03) -- (VOUT_NODE);
% El colector roba corriente de base al transistor de paso
\draw[signal] (Qlim.C) |- (30.2, 6.0) node[circ, fill=black, inner sep=1.2pt] {};

% --- Riel de Salida Regulada (V_OUT) ---
\draw[rail] (VOUT_NODE) -- (41.2, 0 |- Qpass.E);
\node[font=\sffamily\bfseries\scriptsize, above=2pt] at (36.0, 0 |- Qpass.E) {$V_{\text{OUT}}$};

% --- Divisor de Retroalimentación Resistivo (R1, R2) ---
\coordinate (FB_TAP) at (38.0, 4.0);
\draw[signal] (38.0, 0 |- Qpass.E) node[circ, fill=black, inner sep=1.2pt] {} to[R, l=$R_1$] (FB_TAP) node[circ, fill=black, inner sep=1.2pt] {} to[R, l=$R_2$, -*] (38.0, 0.0);
\node[font=\sffamily\tiny\bfseries, right=1pt] at (FB_TAP) {$\text{FB}$};

% Lazo de Retroalimentación por Corredor Superior Despejado (Y = 10.0) (Sin texto)
\draw[feedback_bus] (FB_TAP) -- (38.6, 4.0) -- (38.6, 10.0) -- (21.0, 10.0) |- (Q2.B);

% --- Desacoplo de Salida y Carga Final ---
% Capacitor Cerámico de Salida
\draw[signal] (36.6, 0 |- Qpass.E) to[C, l=$C_{\text{out}}$, *-*] (36.6, 0.0);
% Resistencia de Carga Nominal
\draw[signal] (39.6, 0 |- Qpass.E) to[R, l=$R_{\text{load}}$, *-*] (39.6, 0.0);

% Bornes de Salida
\draw[signal] (41.2, 0 |- Qpass.E) to[short, -o] (41.8, 0 |- Qpass.E) node[right, font=\sffamily\bfseries\scriptsize] {$V_{\text{out},+}$};
\draw[signal] (41.2, 0.0) to[short, -o] (41.8, 0.0) node[right, font=\sffamily\bfseries\scriptsize] {$V_{\text{out},-}$};

\end{circuitikz}
\end{document}
